\documentclass[12pt]{article}
\usepackage[utf8]{inputenc}
\usepackage[margin=0.75in]{geometry}
\usepackage{amsmath}
\usepackage{amssymb}
\usepackage{amsthm}
\usepackage{graphicx} % Required for inserting images
\usepackage{array}

\title{Norm Emergence Report: Resistance}
\author{Aaron Avram}
\date{September 16th 2026}

\begin{document}
\theoremstyle{plain}
\newtheorem{theorem}{Theorem}
\newtheorem{proposition}{Proposition}
\newtheorem{lemma}[theorem]{Lemma}
\newtheorem{corollary}[theorem]{Corollary}
\newtheorem{claim}[theorem]{Claim}

\theoremstyle{definition}
\newtheorem{definition}{Definition}
\newtheorem{example}{Example}
\newtheorem{exercise}[theorem]{Exercise}

\theoremstyle{remark}
\newtheorem{remark}[theorem]{Remark}
\newtheorem{note}[theorem]{Note}

\maketitle

\section{Introduction}
This report assesses the previous students' \textit{Report Revision} against the resistance
model of the \textit{Social Reward Economies} and estimates the work needed to adapt the
current norm emergence code to it. Section 2 states the four implementation requirements
the paper imposes, since every verdict below follows from one of them. Sections 3-5 sort the report's
proposals into parts usable as they are, parts that need refinement, and parts that need to be redone.
Section 6 gives an effort estimate and a sequencing recommendation.

One note on scope: the report proposes an attention allocation mechanism but supplies
only the objective function, with no closed form, no update rule and no statement about what it replaces.
I have worked those out, but have not included any of my derivations which I can supply if necessary: Section 4.4 only contains the verdict.

The general assessment of the report: the statement gossip is a mechanism connecting agents and the key to resistance,
is right. Its recurring error is treating gossip fragmentation, reputational asymmetry and subgroup membership
as structures to be imposed, when the paper treats all three as emergent. Most of the decisions below follow from this observation.

\section{What the Paper Requires of the Implementation}
In the paper \textit{Social Reward Economies}, resistance is not modelled as a product
of generic deviance or systemic disruptions. Instead, it is viewed as the emergence of an alternate reward economy.
Specifically, when a subset $G_r \subseteq G$ of the set of total actors views the dominant schema $\sigma_d$ as detrimental relative to an alternative $\sigma_r$,
its members can form an alternative economy that rewards behaviours aligned with $\sigma_r$ instead of $\sigma_d$.
Concretely, we will view this as a gap in rewards on a subset $S_r$ of the total set of states $S$ where $\sigma_r$ generates significantly better rewards
for $G_r$ than $\sigma_d$. Resistance
is therefore not a property of individual agents but of a newly formed competing reward structure.
The goal of our simulation is to reproduce this formation process from end-to-end and its three possible outcomes.
These are described in the paper as (Eqs. (1)-(2) in section 6.2.2):
\begin{enumerate}
    \item[1.] $G_r$'s weighted feedback dominates $G \setminus G_r$'s for both schema which implies the alternative schema overturns
    the old norm (this is successful resistance / norm change, exemplified by \#MeToo in Section 6.2).
    \item[2.] The two economies' weighted feedback is comparable, this causes polarization: the economies coexist and
    compete without either dominating (Section 7.2.).
    \item[3.] $G \setminus G_r$'s feedback dominates implies $G_r$ fails to change the norm and remains a marginalized subgroup
    (Section 6.2.2, end).
\end{enumerate}
The current engine cannot produce any of the above scenarios.

This is because each gossip update is of the form:
\[ s_i(k, t+1) = \frac{1}{|\mathcal A_p(t)|}\sum_{j \in \mathcal A_p(t)}s_j(k, t) + (v_i(k, t) - v_i(k, t - 1)) \]
where $s_i(k, t)$ is $i$'s reputation estimate for $k$ at time $t$, $\mathcal A_p(t)$ is the set of participating actors
at time $t$ and $v_i(k, t)$ is $i$'s estimate of its personal benefit from $k$ at time $t$.

Now as $t \to \infty$ theoretically we find that
\[ s_i(k, t) \approx s_j(k, t) \]
as the consensus term is identical for all participants and $v_i(k, t)$ plateaus as $t \to \infty$
which implies $v_i(k, t + 1) - v_i(k, t) \to 0$.

The above formula can be verified empirically for sufficiently large $t$ close to norm convergence and it drives
the adoption of the current leader's norm and prevents any resistance from forming.

To augment the current code to test the three outcomes above requires four major changes:
\begin{enumerate}
    \item[1.] \textbf{Persistent personal-utility tracking.} Currently, an agent only acts under its leader's policy,
    and it overwrites its personal-utility estimates whenever it optimizes reputation or status. Thus an agent in $G_r$
    following the dominant norm can never discover the gap $u_i(s, \sigma_r) \gg u_i(s, \sigma_d)$ which is precisely
    what the paper identifies as the trigger for resistance. This is the most important part of the simulation to change.
    \item[2.] \textbf{Sampling $\sigma_r$.} If an agent simply imitates its current leader,
    they will never encounter an alternative schema in the first place so misalignment can be tracked
    but never observed. Thus, followers must have some non-zero probability of deviating from the leader's policy.
    \item[3.] \textbf{Group-relative feedback and reputation.} The quantities $f^r_i$ vs. $f_j$ and $w^r_j$ vs. $w_j$
    described in Section 6.2.1 of the paper are not the same object computed twice, they are different quantities that a subgroup
    computes for itself. Therefore, the implementation needs a mechanism by which a subgroup's internal reputational
    structure can diverge from the group's. Moreover, Section 4.1 of the paper describes gossip networks and opinion leaders
    as emergent structural properties of the reward economy rather than externally imposed ones. Hence the divergence
    in the gossip network should arise from the dynamics of the simulation, not be hard coded as two separate networks
    or two separate parameter sets.
    \item[4.] \textbf{A reward structure encoding misalignment.} The paper's setup gives four
    distinct payoff distributions $(\sigma_r/\sigma_d \times G_r / G \setminus G_r)$, the current reward model
    has no way of representing this and need to be extended.
\end{enumerate}

\section{Usable As They Are}
\begin{itemize}
    \item \textbf{The reward table structure} (Section 5.1 Table 1 in the report). Sampling $u_i(s, \sigma)$
    from group-conditioned normals on $s \in S_r$ with $\text{mean}_r^{G_r} \gg \text{mean}_d^{G_r}$ and $\text{mean}_d^G \gg \text{mean}_r^G$
    is a direct and minimal implementation of the fourth requirement that layers onto the existing reward table construction. This is the only piece
    that needs absolutely no modification.
    \item \textbf{The misalignment condition itself} (Section 3.3 in the report). Defining $G_r$ and $S_r$
    by $u_i(s, \sigma_r) \gg u_i(s, \sigma_d)$ is the right implementation of Section 6.1 of the paper. What needs some reworking
    is the discovery mechanism built on top of it.
    \item \textbf{Diagnostics and hypotheses 1-4} (Section 5.1 in the Report). Subgroup size over time, per-agent and per-group reward totals,
    and the predictions about $|S_r|$, $|G_{\text{unhappy}}|$, the utility gap and within-group variance all
    factor into the size-versus-intensity tradeoff of Section 6.2.3 of the paper and can be carried over directly. The parameters
    they range over need revision (These are addressed in Sections 4.5, 5.4 below).
\end{itemize}

\section{Sensible But Needing Refinement}

\subsection{Similarity-gated gossip}
The report's similarity distance, thresholded at a constant $C$ to decide whether two agents may gossip,
is the piece closest to the paper's own account of how alternative reputational structures form: it makes
the emergence of a $G_r$-specific network and $G_r$-specific weights a testable outcome
rather than a premise, which is the third requirement.

The utility matrix is already computed at each timestep, so implementation only requires accumulating these signals over time
and a threshold check in the gossip algorithm.

Two changes are needed. The report defines the distance over $u_i(s, \sigma_r)$ for pairs already known to be in $G_r$
which is circular (discussed in Section 5.2 below). It should instead be computed over whatever actions were actually observed,
which is well defined before any subgroup exists. And the averaging set, which the report leaves ambiguous,
should be all of S, since $S_r$ is not known in advance.

With this definition the mechanism begins fragmenting the network before anyone deviates, because distances
are computed on utilities derived under $\sigma_d$ and members of $G_r$ already score poorly there. So similarity
and the deviation rate defined below are complementary: similarity separates the network and deviation supplies the actual sample of $\sigma_r$
to members of $G_r$.

\subsection{Discovery of misalignment}
The report does give a defection trigger (Section 3.5 Step 1): agent $i$ optimizes utility in $s$
when
\[ u_i^t(s, \sigma_i) - u_i^t(s, \sigma) > R^{G, t}_i(s, \sigma) - R_i^{G, t}(s, \sigma_i) \]
this presupposes that $i$ already holds an estimate of $u_i(s, \sigma_r)$. It specifies when to defect given the estimate,
never how the estimate is acquired by an agent optimizing reputation and the current code does not supply a solution.

The fix has two parts. Personal utility estimates must be kept current and must not be overridden by leader weights.
And I propose a deviation rate $\mu_{\text{dev}}$, where $\theta(\mu_{\text{dev}})$ is the probability that
an agent optimizes personal utility at a given timestep regardless of its assigned role.

At the moment $\mu_{\text{dev}}$ is fixed, but if it is implemented then a possible extension is to optimize this rate as well.

\subsection{The Dual reputation structure}
The $2\times 2$ reputation matrix of Section 3.2 in the report has the right underlying structure and should be kept as it is essentially
the third requirement, and corresponds to the paper's $w_j$ versus $w^r_j$. Members of $G_r$ maintain two estimates, one
aggregated over $G_r$ and $S_r$ and one over $G$ and $S$; members of $G \setminus G_r$ maintain only the
latter. What should be removed is the four $\gamma$ factors indexing it (Section 5.1 of the report).
 
Once the four $\gamma$'s are removed the report's four-case learning derivation in section 3.4 is largely redundant. Cases 3.4.3 and 3.4.4
are the existing update summed over $G$; case 3.4.1 is the same update summed over $G_r$ and $S_r$. One
update parameterized by its index sets is less code and closer to the paper, which treats the subgroup
economy as the same mechanism over a different population. Case 3.4.2 is separate and discussed in Section 5.3 below.

\subsection{Attention allocation}
The report proposed a model for the rates $\mu_i(s)$ where each agent maximizes $\sum_{s} \theta(\mu_\text{min} + \mu_i(s)) u_i(s,\pi_i(s))$ subject
to a total budget $M$. This is worth keeping, because it implements the second half of Section 6.2.3's success
condition in the paper. Namely, that resistance requires a disproportionately strong and coherent feedback from
$G_r$ on top of a utility gap. If deviating agents are allowed to concentrate their budget on $S_r$,  the feedback loop will emerge from the dynamics of the simulation.
 
I have derived a closed form for the optimum and a $\mathcal O(|G||S|\log|S|)$ method for computing it, together with an averaged update rule
satisfying the usual stepsize conditions, and I propose that this replaces the current $\mu_{i,a}$ with the
same treatment available for $\mu_{i,p}$.

It is important to note that I solved the uniform-$\mu_\text{min}$ case
first for tractability, but section 4.2.2 of the report assumes $\mu_\text{min}(s) \approx 0$ on $S_r$ and larger
elsewhere, and that asymmetry is what drives its predicted split. Uniform $\mu_\text{min}$ weakens the
concentration effect, so a situation-dependent version can be added after the uniform case is implemented.

\subsection{Simulation design}
Varying one parameter per run assumes additive effects. Section 6.2.3 of the paper describes resistance as a joint function
of $|G_r|$ and feedback intensity, so simulations only varying a single-parameter cannot test the paper's success condition at
all. A grid over even two or three parameters such as  $|G_r|$, the reward gap $m_r^{G_r} - m_d^{G_r}$, and now
$\mu_\text{dev}$ and the similarity threshold $C$, is more informative than any number of one-dimensional
runs.
 
\subsection{The concurrent-situations experiment}
Section 5.2 of the report suggests varying $\mu_\text{min}(s)$ to test institutional participation floors which is worth keeping and
should run before even testing resistance as it validates the attention mechanism in the single-economy
setting we already understand. Additionally, debugging the optimizer and two-economy dynamics simultaneously could cause
silent failure. Additionally, the report's hypotheses about the effect of the rate on resistance should be restated in terms of the optimization of the rate proposed in Section 4.4 above.

\section{Needing to Be Redone}
 
\subsection{Manually fragmented networks and separate $\gamma$'s}
The report has agents in $G_r$ participate in two gossip networks and agents in $G\setminus G_r$ in one,
with membership determining connectivity. But membership is what the simulation is supposed to produce, so
connectivity cannot be dependent on it. Section 4.1 of the paper describes gossip networks as capable of
both reinforcing and destabilizing power structures through ongoing reputational dynamics, and Section 7.2.1
describes polarized subgroups as becoming insulated through their own feedback amplification. Therefore insulation
is an outcome, not an input. The similarity threshold of Section 4.1 above replaces this.
 
The same objection applies to the four $\gamma$'s (and the four $\kappa$'s from the October report), which
index reputation weight by group membership to manufacture the asymmetry that should emerge from
fragmentation. Section 6.2.1 in the paper is explicit that agents in $G_r$ acquire higher $w^r_j$ by supporting
$\sigma_r$ and condemning the dominant schema. Therefore, $w^r_j$ should be an output of the feedback process, not a
parameter. This also makes hypotheses 5--7 of Section 5.1 in the report unanswerable in any interesting sense, since they
would test the consequences of an assumption rather than of a mechanism.

\subsection{The subgroup emergence mechanism}
Section 3.5 of the Report has a four-step account which is circular at two levels. Step 1's trigger requires an estimate the agent
cannot have (Section 4.2 above). Step 2 forms a utility gossip link only when both agents are in a diverging
situation $s \in S(\sigma_i,t) \cap S(\sigma_j,t)$ and have aligned on $\sigma_i = \sigma_j =
\sigma_r$, but $S(\sigma_r) = S_r$ and $G(\sigma_r) = G_r$ by the report's own definitions, so the
conditions for forming the links that produce $G_r$ presuppose $G_r$ and $S_r$. Step 3's transitive
expansion then propagates an assumed property.
 
The replacement is clear, the first requirement above gives agents the estimates the trigger needs, $\mu_\text{dev}$
gives them the samples, and the similarity threshold computed over observed actions does the work of the
utility gossip links without knowing who is in $G_r$. The report's distinction between utility gossip links
and reputation gossip links collapses into a single similarity-gated link, which is one fewer structure to
build.

\subsection{Out-group targeting}
Section 3.4.2 of the report defines a mechanism for identifying the lowest reputation non-member, with its own tie
threshold, and then never uses it. Nothing in Section 5's simulation plan consumes $I_j(t)$. It has no
counterpart in the paper, which models the out-group relation through relative feedback weight, not
individual targeting. I would cut this part instead of implementing it with no usecase.
 
\subsection{Seeded initialization}
The report seeds $G_r$ with $\sigma_r$ at $t = 0$. This answers a survival question, does an existing
alternative economy persist? This is not the formation question that both the report and Section 6.1 are
about. Seeding was previously unavoidable, since without requirements 1 and 2 from above nothing could form; with them in place,
every agent should start under $\sigma_d$ and $G_r$ should be an outcome we measure rather than a condition
we impose.

\section{Effort Estimate and Sequencing}
 
I divided the changes into three tiers based on how difficult it will be to integrate them into the current codebase.
The risk is concentrated in the code for tracking metrics rather than in any
of the new mechanisms.
 
\textbf{Low.} (1) The reward table extension: this layers onto the existing reward table structure. (2) Similarity score: the
utility matrix already exists, and tracking similarity can be added independently without affecting any other part of the code. (3) $\mu_\text{dev}$: one sampling
step in role selection.
 
\textbf{Medium.} (4) Persistent personal-utility tracking: small in code volume but it touches the core
agent update loop, specifically the path where a reputation-optimizing agent copies its leader's weights.
Everything downstream depends on it, so it should be done first and verified against current single-norm
results. (5) Dual reputation bookkeeping: this requires parameterizing the existing update by index sets to compare reputations
between the groups.
 
\textbf{High.} (6) The attention mechanism: the constrained optimizer, implementing the solution for the closed form and the averaged
update, then the same for $\mu_{i,p}$. This is the largest piece of new code, and it changes a quantity many
other parts read. (7) Updating the metrics layer to accept multiple leaders. This is a major risk because 
it fails silently. \texttt{paper\_welfare} takes a single \texttt{leader\_id}; the
leader metric takes a global argmax of follower counts; convergence detection and the plotting harnesses
inherit the same assumption. With two norms present these keep producing plausible numbers that mean
nothing. Every metric needs auditing before any two-economy result is trusted.
 
\textbf{Ordering.} Items 4 and 3 must be implemented before all of the others: until personal utility is tracked and followers can
deviate, no configuration of the rest can produce a second economy. Item 7 also needs to be handled early to avoid any incorrect data
on actual experiments. A reasonable order is
4, 3, 1, 2, 7, then 6 validated on the single-economy experiment mentioned in section 4.6 above, and finally 5 with the full
two-economy grid discussed in section 4.5.

\end{document}